\section{Related Works}
% @Ziyan @Rui @Can
% \begin{itemize}
%     \item Multimodal embedding benchmarks.
%     \item Video Representation learning.
%     \item Visual Document representation learning.
% \end{itemize}
\subsection{Multimodal Embedding Benchmarks}
Numerous benchmarks have been proposed to evaluate multimodal models, with most early efforts focusing on static image-text pairs. Datasets such as MSCOCO~\citep{mscoco}, Flickr30K~\citep{flickr30k}, and Conceptual Captions~\citep{ccdataset} enabled progress in tasks like image captioning and retrieval. Linear probing is also a common evaluation setting that trains a linear layer for image classification~\citep{clip} to investigate the generalization of representation vectors. More recent benchmarks like M-BEIR~\citep{uniir} and MMEB~\citep{vlm2vec} introduced multi-task evaluations for multimodal embedding models, covering tasks such as retrieval and QA. However, these benchmarks remain limited to static images and short contexts. Video-based benchmarks such as MSR-VTT~\citep{msrvtt}, QVHighlights~\citep{qvhighlights}, and ActivityNet Captions~\citep{activitynet} target retrieval and captioning tasks, yet lack unified evaluation frameworks for embeddings. Our \benchmark addresses these gaps by providing a comprehensive embedding benchmark for diverse modalities.


\subsection{Video Representation Learning}
Video representation learning has evolved significantly, progressing from early convolutional approaches to sophisticated transformer-based architectures. Traditional vision-language models like CLIP~\citep{clip} and BLIP~\citep{blip}, while effective for image-text tasks, often struggle to capture the temporal dynamics inherent in video data. To address this, recent models have been developed to better handle the complexities of video understanding. VideoCLIP~\citep{videoclip}, VideoCoCa~\citep{videococa} integrates contrastive learning with captioning objectives to enhance video-text representation alignment. InternVideo2~\citep{internvideo2} employs a progressive training approach that unifies masked video modeling, cross-modal contrastive learning, and next-token prediction, resulting in superior performance on over 60 video and audio tasks. 
% VideoPrism~\citep{videoprism} serves as a foundational visual encoder designed for general-purpose video understanding, achieving state-of-the-art results across a wide spectrum of tasks. 
Recent models like LLaVE~\citep{llave} and  LamRA~\citep{lamra}, though trained exclusively on image-text data, have demonstrated the ability to generalize to text-video retrieval tasks in a zero-shot manner. These advancements highlight the ongoing efforts to develop models capable of effectively understanding and representing the complex temporal and semantic information in video data.



\subsection{Visual Document Representation Learning}

Visual document representation learning has become increasingly vital for tasks such as document retrieval and retrieval-augmented generation. Traditional text-based models often struggle to capture the rich visual and structural information present in documents, necessitating approaches that integrate both visual and textual modalities.
One notable advancement is ColPali~\citep{colpali}, which leverages vision-language models to enhance document retrieval efficiency by effectively capturing both textual and visual features. In the realm of retrieval-augmented generation, VisRAG~\citep{visrag} establishes a vision-based RAG pipeline that directly embeds documents as images using vision-language models, thereby preserving the original document information and outperforming traditional text-based RAG systems. Similarly, ViDoRAG~\citep{vidorag} introduces a multi-agent framework tailored for complex reasoning across visual documents, employing a dynamic iterative reasoning process to enhance retrieval and generation tasks. 
Furthermore, benchmarks like MMLongBench-Doc~\citep{mmlongbench} have been developed to assess long-context document understanding with visualizations, providing a comprehensive evaluation framework for multimodal models. 

\subsection{Unified Modality Representation Learning}

Unified modality representation learning aims to build a single model capable of processing information across diverse data types, such as text, images, audio, and video. Some methods, like Uni-Retrieval~\citep{jia2025uni}, leverage multimodal large language models (MLLMs) and prompt-tuning to accommodate a variety of queries, achieving strong performance on universal benchmarks. While other recent studies have achieved remarkable results on diverse text and image tasks~\citep{gme,b3,metaembed,vlm2vec}, these models were not designed to unify image, video, and visual document retrieval within a single framework. Our latest work, \model, is the first to address this specific challenge. Following its release, subsequent efforts such as TTE~\citep{tte} and Seed1.6-Embedding~\citep{seed16embedding} have proposed advanced techniques and demonstrated strong performance on our \benchmark benchmark.

% Meanwhile, methods such as UniversalRAG~\citep{yeo2025universalrag} and UniRAG~\citep{sharifymoghaddam2025unirag} improve retrieval-augmented generation by dynamically routing queries to the most suitable modality and granularity, enhancing both flexibility and accuracy. 
